\section{Conclusion}

This thesis studied the problem of constructing a post-quantum threshold batch identity-based encryption scheme. The motivation comes from encrypted mempool applications, where transaction contents should remain hidden until their ordering has been fixed, while the decryption authority should be distributed among multiple parties to avoid single-point trust. Existing batch threshold encryption schemes are often based on pairing assumptions, which are not secure against quantum adversaries. Therefore, developing a lattice-based alternative is an important direction toward post-quantum secure encrypted mempools.

To address this problem, this thesis combines lattice-based batched decryption techniques with partial lattice trapdoors. In the proposed construction, the decryption capability of a lattice trapdoor is split among multiple authorities. Each party only holds a partial trapdoor and can generate a partial pre-decryption key locally, while only a qualified threshold set can combine their shares to recover the final batch decryption key. This gives a direct algebraic approach for thresholdizing lattice-based batched identity-based encryption.

The thesis also analyzes the correctness and security of the proposed scheme. Correctness follows from the structure of the shifted multi-preimage sampler, the explainable SampleLeft procedure, and the reconstruction property of partial lattice trapdoors. For security, the construction is analyzed in a highly selective security model through a sequence of hybrid games, relying on the hardness of lattice assumptions such as LWE and SIS, together with statistical indistinguishability arguments from Gaussian sampling and noise flooding.

Overall, the proposed scheme provides a theoretical step toward post-quantum threshold encryption for encrypted mempools. Although several parts of the construction still require further refinement, especially in parameter selection, implementation, and tighter security analysis, the work demonstrates that partial lattice trapdoors can serve as a promising tool for distributing decryption authority in lattice-based identity-based encryption systems.